% titel: Blatt 1, A2: Monoisotopische Elemente / A7: Isotope
% Teil II: der durchgerechnete Weg.

\subsection*{A2 --- Monoisotopische Elemente}
Recherche-Aufgabe (ChatGPT vs.\ Wikipedia): Ein \textbf{monoisotopisches
Element} hat nur ein stabiles Isotop. Beispiele: Beryllium ${}^{9}$Be,
Fluor ${}^{19}$F, Natrium ${}^{23}$Na, Aluminium ${}^{27}$Al, Phosphor
${}^{31}$P, Gold ${}^{197}$Au. Die Antwort von KI und Wikipedia sollten hier
übereinstimmen; Abweichungen (z. B.\ KI nennt fälschlich ein Element mit
sehr dominantem Isotop) als Fehler der KI erkennen.

\subsection*{A7 a) Warum keine Isotope?}
Die zwei Kerne haben \textbf{verschiedene Ordnungszahlen}: Sie gehören zu
zwei verschiedenen Elementen (B und C). Isotope haben gleiche $Z$ und nur
verschiedene $N$. Gleiche Massenzahl, aber andere $Z$, nennt man
\textbf{Isobare} --- keine Isotope.

\subsection*{A7 b) Isotopenhäufigkeit; ein Reinelement}
\textbf{Isotopenhäufigkeit:} der prozentuale Anteil eines Isotops an der
Gesamtzahl der Atome des Elements in der Natur. Beispiel: Ein \textbf{Reinelement}
wie Fluor hat nur ein stabiles Isotop, Häufigkeit $\approx 100\%$.

\begin{ergebnis}
Zwei Kerne mit verschiedener $Z$ sind keine Isotope (sie sind Isobare).
Isotopenhäufigkeit = relativer Anteil eines Isotops; Reinelement z. B.\ gefunden
in A2 (Beryllium, Fluor, \dots{} ).
\end{ergebnis}

\begin{kontrolle}
Die Bedingung ``gleiche $Z$ = Isotope" ist hier verletzt, also korrekt
ausgeschlossen. Die gemischte Nomenklatur (Isotop/Isobar) typische Falle.
\end{kontrolle}

\begin{merksatz}
Isotope: gleiche $Z$, verschiedene $A$. Isobare: gleiche $A$ ($N+Z$),
verschiedene $Z$. Das $Z$ entscheidet über das Element.
\end{merksatz}
